%% ==================================================================
%% Consolidated bibliography
%% ==================================================================

\begin{thebibliography}{99}

\bibitem{meredith2005rho}
L.G.~Meredith and M.~Radestock. \newblock A reflective higher-order calculus. \newblock In \emph{Proceedings of ETAPS/FOSSACS}, 2005.

\bibitem{leifer2000deriving}
J.J.~Leifer and R.~Milner. \newblock Deriving bisimulation congruences for reactive systems. \newblock In \emph{Proceedings of CONCUR 2000}, LNCS 1877, pp.~243--258, 2000.

\bibitem{gillespie1977exact}
D.T.~Gillespie. \newblock Exact stochastic simulation of coupled chemical reactions. \newblock \emph{Journal of Physical Chemistry}, 81(25):2340--2361, 1977.

\bibitem{cao2021quantum}
Y.~Cao, J.~Lu, and Y.~Lu. \newblock Quantum Gillespie algorithm and related methods for open quantum systems. \newblock Preprint, 2021. \newblock (See also: Breuer \& Petruccione, \emph{The Theory of Open Quantum Systems}, OUP.)

\bibitem{danos2004reversible}
V.~Danos and J.~Krivine. \newblock Reversible communicating systems. \newblock In \emph{Proceedings of CONCUR 2004}, LNCS 3170, pp.~292--307, 2004.

\bibitem{bombelli1987space}
L.~Bombelli, J.~Lee, D.~Meyer, and R.D.~Sorkin. \newblock Space-time as a causal set. \newblock \emph{Physical Review Letters}, 59(5):521--524, 1987. % ---- Part II references ----

\bibitem{abramsky1991domain}
S.~Abramsky. \newblock Domain theory in logical form. \newblock \emph{Annals of Pure and Applied Logic}, 51(1--2):1--77, 1991.

\bibitem{fagin1995reasoning}
R.~Fagin, J.Y.~Halpern, Y.~Moses, and M.Y.~Vardi. \newblock \emph{Reasoning About Knowledge}. \newblock MIT Press, 1995.

\bibitem{hatcher2002algebraic}
A.~Hatcher. \newblock \emph{Algebraic Topology}. \newblock Cambridge University Press, 2002.

\bibitem{milner1989communication}
R.~Milner. \newblock \emph{Communication and Concurrency}. \newblock Prentice Hall, 1989.

\bibitem{wilson1974renormalization}
K.G.~Wilson and J.~Kogut. \newblock The renormalization group and the $\epsilon$ expansion. \newblock \emph{Physics Reports}, 12(2):75--199, 1974. % ---- Part III references ----

\bibitem{turing1939}
A.M.~Turing. \newblock Systems of logic based on ordinals. \newblock \emph{Proceedings of the London Mathematical Society}, 45(1):161--228, 1939.

\bibitem{sutton2018reinforcement}
R.S.~Sutton and A.G.~Barto. \newblock \emph{Reinforcement Learning: An Introduction}. \newblock 2nd edition. MIT Press, 2018.

\bibitem{friston2010free}
K.~Friston. \newblock The free-energy principle: a unified brain theory? \newblock \emph{Nature Reviews Neuroscience}, 11(2):127--138, 2010.

\bibitem{tononi2008consciousness}
G.~Tononi. \newblock Consciousness as integrated information: a provisional manifesto. \newblock \emph{Biological Bulletin}, 215(3):216--242, 2008. % ---- Origins of Life references ----

\bibitem{cronin2023assembly}
L.~Cronin and S.~Walker. \newblock Assembly theory explains and quantifies selection and evolution. \newblock \emph{Nature}, 622:321--328, 2023.

\bibitem{smith2016origin}
E.~Smith and H.J.~Morowitz. \newblock \emph{The Origin and Nature of Life on Earth: The Emergence of the Fourth Geosphere}. \newblock Cambridge University Press, 2016.

\bibitem{regev2001representation}
A.~Regev, W.~Silverman, and E.~Shapiro. \newblock Representation and simulation of biochemical processes using the $\pi$-calculus process algebra. \newblock \emph{Pacific Symposium on Biocomputing}, 6:459--470, 2001.

\bibitem{KnotsNote}
L.\ G.\ Meredith. \emph{Knots as cost-accounted rholang processes: a Dowker--Thistlethwaite-driven encoding and weak-bisimulation invariance.} F1R3FLY.io, 2026.

\bibitem{ClassifierNote}
L.\ G.\ Meredith. \emph{Classifying knot diagrams with a spatial--behavioural logic: a logic obtained for free from the encoding of knots as processes.} F1R3FLY.io, 2026.

\bibitem{Omnibus}
\emph{Graph-structured lambda theories: a reference for the working developer.} F1R3FLY.io, 2026.

\bibitem{Probing}
L.\ G.\ Meredith. \emph{Probing universality: context-labelled bisimulation and a relative refinement of Turing completeness for interactive GSLTs.} F1R3FLY.io, 2026.

\bibitem{CostMonad}
L.\ G.\ Meredith. \emph{Cost accounting as a monad for continued graph-structured lambda theories.} F1R3FLY.io, 2026.

\bibitem{History}
L.\ G.\ Meredith. \emph{History as an endofunctor.} F1R3FLY.io, 2026.

\bibitem{StochSim}
L.\ G.\ Meredith. \emph{Stochastic simulation for decorated ciGSLTs.} F1R3FLY.io, 2026.

\bibitem{Zeus}
L.\ G.\ Meredith. \emph{Splitting the mind of Zeus: realizability, collection monads, and the choice content of the McBride derivative in a finite-limits metalanguage.} F1R3FLY.io, 2026.

\bibitem{NamespaceLogic}
L.\ G.\ Meredith and M.\ Radestock. Namespace logic: a logic for a reflective higher-order calculus. In \emph{Trustworthy Global Computing}, LNCS 3705, Springer, 2005.

\bibitem{RHO}
L.\ G.\ Meredith and M.\ Radestock. A reflective higher-order calculus. \emph{Electronic Notes in Theoretical Computer Science} 141(5):49--67, 2005.

\bibitem{OSLF}
\emph{Operational semantics in logical form (OSLF): a family of algorithms auto-generating logics and type systems from graph-structured lambda theories.}

\bibitem{Cleaveland}
R.\ Cleaveland. On automatically explaining bisimulation inequivalence. In \emph{Computer Aided Verification (CAV)}, LNCS 531, Springer, 1990.

\bibitem{Caires}
L.\ Caires. Behavioral and spatial observations in a logic for the pi-calculus. In \emph{Foundations of Software Science and Computation Structures (FoSSaCS)}, LNCS 2987, Springer, 2004.

\bibitem{CairesCardelli}
L.\ Caires and L.\ Cardelli. A spatial logic for concurrency (part I). \emph{Information and Computation} 186(2):194--235, 2003.

\bibitem{Milner}
R.\ Milner. Functions as processes. \emph{Mathematical Structures in Computer Science} 2(2):119--141, 1992.

\bibitem{Plotkin}
G.\ D.\ Plotkin. A structural approach to operational semantics. \emph{Journal of Logic and Algebraic Programming} 60--61:17--139, 2004.

\bibitem{Wigner}
E.\ P.\ Wigner. The unreasonable effectiveness of mathematics in the natural sciences. \emph{Communications on Pure and Applied Mathematics} 13(1):1--14, 1960.

\bibitem{Hamming}
R.\ W.\ Hamming. The unreasonable effectiveness of mathematics. \emph{The American Mathematical Monthly} 87(2):81--90, 1980.

\bibitem{Landauer}
R.\ Landauer. Irreversibility and heat generation in the computing process. \emph{IBM Journal of Research and Development} 5(3):183--191, 1961.

\bibitem{Hinton}
G.\ Hinton. The forward-forward algorithm: some preliminary investigations, 2022. (Mortal computation.)

\bibitem{Noble}
D.\ Noble. \emph{Dance to the Tune of Life: Biological Relativity.} Cambridge University Press, 2016.

\bibitem{SmithMorowitz}
E.\ Smith and H.\ J.\ Morowitz. \emph{The Origin and Nature of Life on Earth: The Emergence of the Fourth Geosphere.} Cambridge University Press, 2016.

\bibitem{Gillespie}
D.\ T.\ Gillespie. Exact stochastic simulation of coupled chemical reactions. \emph{Journal of Physical Chemistry} 81(25):2340--2361, 1977.

\bibitem{Newell}
A. Newell and H. A. Simon. \emph{Human Problem Solving.} Prentice-Hall, 1972.

\bibitem{Bouton}
C. L. Bouton. Nim, a game with a complete mathematical theory. \emph{Annals of Mathematics} 3(1/4):35--39, 1901--1902.

\bibitem{ONAG}
J. H. Conway. \emph{On Numbers and Games.} Academic Press, 1976.

\bibitem{Hardin}
G. Hardin. The competitive exclusion principle. \emph{Science} 131(3409):1292--1297, 1960.

\bibitem{Tilman}
D. Tilman. \emph{Resource Competition and Community Structure.} Princeton University Press, 1982.

\bibitem{Margulis}
L. Margulis. \emph{Origin of Eukaryotic Cells.} Yale University Press, 1970.

\bibitem{compaut}
F1R3FLY.io, \emph{Compression versus Autonomy: an error-theoretic account of an architectural crossover common to biology, silicon, and distributed systems}, research note, 2026-05-25.

\bibitem{CARho}
F1R3FLY.io, \emph{Cost-accounted rho calculus}, research note, 2026.

\bibitem{CE}
F1R3FLY.io, \emph{The cost endofunctor on continued interactive GSLTs}, research note, 2026.

\bibitem{vtok}
F1R3FLY.io, \emph{The Virtual Token: no-arbitrage, conversion, and the cost-accounted transport of energy}, research note, 2026.

\bibitem{finrho}
F1R3FLY.io, \emph{Financial contracts in the cost-accounted rho calculus}, research note, 2026.

\bibitem{MC}
F1R3FLY.io, \emph{Mortal computation}, research note, 2026.

\bibitem{SPJ2000}
S.~Peyton~Jones, J.-M.~Eber, J.~Seward, \emph{Composing contracts: an adventure in financial engineering}, ICFP 2000.

\bibitem{foster}
R.~M.~Foster, \emph{The average impedance of an electrical network}, in \emph{Contributions to Applied Mechanics}, 1949, pp.~333--340.

\bibitem{hoory}
S.~Hoory, N.~Linial, A.~Wigderson, \emph{Expander graphs and their applications}, Bull.\ AMS 43(4):439--561, 2006.

\bibitem{lindeman}
R.~L.~Lindeman, \emph{The trophic-dynamic aspect of ecology}, Ecology 23(4):399--417, 1942.

\bibitem{noe}
R.~No\"e, P.~Hammerstein, \emph{Biological markets: supply and demand determine the effect of partner choice in cooperation, mutualism and mating}, Behavioral Ecology and Sociobiology 35:1--11, 1994.

\bibitem{kiers}
E.~T.~Kiers et al., \emph{Reciprocal rewards stabilize cooperation in the mycorrhizal symbiosis}, Science 333:880--882, 2011.

\bibitem{west}
G.~B.~West, J.~H.~Brown, B.~J.~Enquist, \emph{A general model for the origin of allometric scaling laws in biology}, Science 276:122--126, 1997.

\bibitem{coase}
R.~H.~Coase, \emph{The nature of the firm}, Economica 4(16):386--405, 1937.

\bibitem{hanson}
R.~Hanson, D.~Martin, C.~McCarter, J.~Paulson, \emph{If loud aliens explain human earliness, quiet aliens are also rare}, The Astrophysical Journal 922:182, 2021.


\bibitem{qcs}
L.~G. Meredith. \emph{Quoting is Colour-Swap: a model of the rho calculus in the
category of bags}. F1R3FLY.io research note, 2026.

\bibitem{paths}
L.~G. Meredith. \emph{Paths are Subspaces: a cut-triggered polymorphism of RSpace}.
F1R3FLY.io research note, 2026.

\bibitem{choicesched}
L.~G. Meredith. \emph{From the Axiom of Choice to Fair Scheduling: toward a
Curry--Howard correspondence for nondeterminism, in which choice principles are the
types and fair schedulers are the terms}. F1R3FLY.io research note, June 2026.

\bibitem{knot}
L.~G. Meredith. \emph{A Knotted Universe: a new notion of reflective set theories}.
F1R3FLY.io research note, 2026.

\bibitem{weihrauch}
V. Brattka, G. Gherardi, and A. Pauly. Weihrauch complexity in computable analysis.
In \emph{Handbook of Computability and Complexity in Analysis}, Springer, 2021.

\bibitem{soare}
R.~I. Soare. \emph{Turing Computability: Theory and Applications}. Springer, 2016.

\bibitem{copeland}
B.~J. Copeland. Hypercomputation. \emph{Minds and Machines}, 12(4):461--502, 2002.

\bibitem{leifermilner}
J.~J. Leifer and R. Milner. Deriving bisimulation congruences for reactive systems.
In \emph{CONCUR 2000}, LNCS 1877, pp. 243--258, 2000.

\bibitem{sewell}
P. Sewell. From rewrite rules to bisimulation congruences.
\emph{Theoretical Computer Science}, 274(1--2):183--230, 2002.

\bibitem{milnerCC}
R. Milner. \emph{Communication and Concurrency}. Prentice Hall, 1989.

\bibitem{rutten}
J.~J.~M.~M. Rutten. Universal coalgebra: a theory of systems.
\emph{Theoretical Computer Science}, 249(1):3--80, 2000.

\bibitem{turiplotkin}
D. Turi and G. Plotkin. Towards a mathematical operational semantics.
In \emph{LICS 1997}, pp. 280--291, 1997.

\bibitem{aczel}
P. Aczel. \emph{Non-Well-Founded Sets}, CSLI Lecture Notes 14. CSLI Publications, 1988.

\bibitem{bell}
J.~S. Bell. On the Einstein Podolsky Rosen paradox.
\emph{Physics Physique Fizika}, 1(3):195--200, 1964.

\bibitem{hensen}
B. Hensen et al. Loophole-free Bell inequality violation using electron spins separated
by 1.3 kilometres. \emph{Nature}, 526:682--686, 2015.

\bibitem{brunner}
N. Brunner, D. Cavalcanti, S. Pironio, V. Scarani, and S. Wehner. Bell nonlocality.
\emph{Reviews of Modern Physics}, 86(2):419--478, 2014.

\bibitem{abramsky}
S. Abramsky and A. Brandenburger. The sheaf-theoretic structure of non-locality and
contextuality. \emph{New Journal of Physics}, 13:113036, 2011.


\bibitem{consensus}
L.~G. Meredith. \emph{From Coalition Logic to Rho Calculus: toward a Curry--Howard
correspondence for consensus protocol synthesis}. F1R3FLY.io research note, 2026.

\bibitem{martinlof}
P. Martin-L\"of. \emph{Intuitionistic Type Theory}. Bibliopolis, Naples, 1984.

\bibitem{spector}
C. Spector. Provably recursive functionals of analysis. In \emph{Recursive Function
Theory}, Proc. Symp. Pure Math. V, pp. 1--27. AMS, 1962.

\bibitem{bbc}
S. Berardi, M. Bezem, and T. Coquand. On the computational content of the axiom of
choice. \emph{Journal of Symbolic Logic}, 63(2):600--622, 1998.

\bibitem{krivine}
J.-L. Krivine. Dependent choice, `quote' and the clock.
\emph{Theoretical Computer Science}, 308(1--3):259--276, 2003.

\bibitem{escardo}
M.~H. Escard\'o and P. Oliva. Selection functions, bar recursion and backward
induction. \emph{Mathematical Structures in Computer Science}, 20(2):127--168, 2010.

\bibitem{francez}
N. Francez. \emph{Fairness}. Texts and Monographs in Computer Science. Springer, 1986.

\bibitem{abramskyPaP}
S. Abramsky. Proofs as processes. \emph{Theoretical Computer Science},
135(1):5--9, 1994.

\bibitem{bellinscott}
G. Bellin and P.~J. Scott. On the $\pi$-calculus and linear logic.
\emph{Theoretical Computer Science}, 135(1):11--65, 1994.

\bibitem{cairespfenning}
L. Caires and F. Pfenning. Session types as intuitionistic linear propositions.
In \emph{CONCUR 2010}, LNCS 6269, pp. 222--236. Springer, 2010.

\bibitem{wadler}
P. Wadler. Propositions as sessions. \emph{Journal of Functional Programming},
24(2--3):384--418, 2014.

\bibitem{diaconescu}
R. Diaconescu. Axiom of choice and complementation.
\emph{Proceedings of the AMS}, 51:176--178, 1975.


\bibitem{Noether}
E.~Noether, \emph{Invariante Variationsprobleme}, Nachrichten von der Gesellschaft der Wissenschaften zu G\"ottingen, Mathematisch-Physikalische Klasse, 1918, pp.~235--257.

\bibitem{Hodge}
L.-H.~Lim, \emph{Hodge Laplacians on graphs}, SIAM Review 62(3), 2020, pp.~685--715.

\bibitem{Bennett}
C.~H.~Bennett, \emph{The thermodynamics of computation --- a review}, International Journal of Theoretical Physics 21(12), 1982, pp.~905--940.

\bibitem{Miller2006}
M.~S.~Miller, \emph{Robust Composition: Towards a Unified Approach to Access Control and Concurrency Control}, PhD thesis, Johns Hopkins University, 2006.

\bibitem{TwoErasures}
F1R3FLY.io, \emph{Two erasures and two laxities: cost, history, and internalisation in continued interactive GSLTs}, research note, 2026.


\bibitem{Kant1783}
I.~Kant, \emph{Prolegomena to Any Future Metaphysics That Will Be Able to Come Forward as Science}, 1783. Translated by G.~Hatfield, Cambridge University Press, revised edition, 2004.

\bibitem{Kant1781}
I.~Kant, \emph{Critique of Pure Reason}, 1781 (A) / 1787 (B). Translated by P.~Guyer and A.~W.~Wood, Cambridge University Press, 1998.

\bibitem{Huet97}
G.~Huet. The zipper. \emph{Journal of Functional Programming}, 7(5):549--554, 1997.

\bibitem{McBride08}
C.~McBride. Clowns to the left of me, jokers to the right: dissecting data structures. In \emph{Proceedings of POPL 2008}, pp.~287--295, 2008.

\bibitem{spacecalc}
L.G.~Meredith and M.~Radestock, \emph{A calculus from outer space}, research note, F1R3FLY.io.



%% ---- imported with the GSLT reference chapters (draft 7) ----

\bibitem{AbramskyGayNagarajan1996}
S.~Abramsky, S.~Gay, and R.~Nagarajan.
\newblock Interaction categories and the foundations of typed concurrent programming.
\newblock In \emph{Deductive Program Design}, Springer, 1996.

\bibitem{BaezWilliams2020}
J.~Baez and C.~Williams.
\newblock Enriched Lawvere theories for operational semantics.
\newblock \emph{Electronic Proceedings in Theoretical Computer Science}, 2020.

\bibitem{Barendregt1991}
H.~Barendregt.
\newblock Introduction to generalized type systems.
\newblock \emph{Journal of Functional Programming}, 1(2):125--154, 1991.

\bibitem{CardelliGordon2000}
L.~Cardelli and A.~D.~Gordon.
\newblock Mobile ambients.
\newblock \emph{Theoretical Computer Science}, 240(1):177--213, 2000.

\bibitem{deBruijn1972}
N.~G.~de~Bruijn.
\newblock Lambda calculus notation with nameless dummies.
\newblock \emph{Indagationes Mathematicae}, 34:381--392, 1972.

\bibitem{kozen}
D.~Kozen. \newblock Results on the propositional $\mu$-calculus. \newblock \emph{Theoretical Computer Science}, 27(3):333--354, 1983.

\bibitem{sangiorgi}
D.~Sangiorgi and D.~Walker. \newblock \emph{The Pi-Calculus: A Theory of Mobile Processes}. \newblock Cambridge University Press, 2001.

\bibitem{HennessyMilner1985}
M.~Hennessy and R.~Milner.
\newblock Algebraic laws for nondeterminism and concurrency.
\newblock \emph{Journal of the ACM}, 32(1):137--161, 1985.

\bibitem{LambekScott1986}
J.~Lambek and P.~J.~Scott.
\newblock \emph{Introduction to Higher Order Categorical Logic}.
\newblock Cambridge University Press, 1986.

\bibitem{Lawvere1963}
F.~W.~Lawvere.
\newblock Functorial semantics of algebraic theories.
\newblock PhD thesis, Columbia University, 1963.

\bibitem{MillerYeeShapiro2003}
M.~S.~Miller, K.-P.~Yee, and J.~Shapiro.
\newblock Capability myths demolished.
\newblock Technical report, Johns Hopkins University, 2003.

\bibitem{Milner1992}
R.~Milner.
\newblock Functions as processes.
\newblock \emph{Mathematical Structures in Computer Science}, 2(2):119--141, 1992.

\bibitem{Plotkin1981}
G.~D.~Plotkin.
\newblock A structural approach to operational semantics.
\newblock Technical Report DAIMI FN-19, Aarhus University, 1981.

\bibitem{Reynolds2002}
J.~C.~Reynolds.
\newblock Separation logic: a logic for shared mutable data structures.
\newblock In \emph{LICS}, 2002.

\bibitem{vanGlabbeek1990}
R.~J.~van~Glabbeek.
\newblock The linear time--branching time spectrum.
\newblock In \emph{CONCUR}, Springer, 1990.

%% ----------------------------------------------------------------
%%  Draft 10: references imported with the Prestige chapters on the
%%  reference bootstrap problem and notationality.
%%  Two keys deduplicate against entries already present:
%%  the source note's meredith2005 is meredith2005rho here, and its
%%  turi1997 is turiplotkin.
%% ----------------------------------------------------------------

\bibitem{abadi2001} M. Abadi and C. Fournet. Mobile values, new names, and secure
communication. \emph{POPL}, 2001.

\bibitem{asai2018} M. Asai and A. Fukunaga. Classical planning in deep latent space:
bridging the subsymbolic-symbolic boundary. \emph{AAAI}, 2018.

\bibitem{asai2019} M. Asai and H. Kajino. Towards stable symbol grounding with
zero-suppressed state autoencoder. \emph{ICAPS}, 2019.

\bibitem{begus2026} G. Begu\v{s}, M. D\k{a}bkowski, R. L. Sprouse, D. F. Gruber and
S. Gero. The phonology of sperm whale coda vowels. \emph{Proceedings of the Royal
Society B} 293(2069):20252994, 2026.

\bibitem{benderkoller2020} E. M. Bender and A. Koller. Climbing towards NLU: on meaning,
form, and understanding in the age of data. \emph{ACL}, 2020.

\bibitem{bengtson2011} J. Bengtson, M. Johansson, J. Parrow and B. Victor. Psi-calculi: a
framework for mobile processes with nominal data and logic. \emph{Logical Methods in
Computer Science} 7(1), 2011.

\bibitem{bisk2020} Y. Bisk, A. Holtzman, J. Thomason, J. Andreas, Y. Bengio, J. Chai,
M. Lapata, A. Lazaridou, J. May, A. Nisnevich, N. Pinto and J. Turian. Experience grounds
language. \emph{EMNLP}, 2020.

\bibitem{castro2021} P. S. Castro, T. Kastner, P. Panangaden and M. Rowland. MICo:
improved representations via sampling-based state similarity for Markov decision
processes. \emph{NeurIPS}, 2021.

\bibitem{estetika2026} On the theory of musical scores and its relation to
work-performance. \emph{Estetika: The European Journal of Aesthetics}, 2026.

\bibitem{ferns2004} N. Ferns, P. Panangaden and D. Precup. Metrics for finite Markov
decision processes. \emph{UAI}, 2004.

\bibitem{gelada2019} C. Gelada, S. Kumar, J. Buckman, O. Nachum and M. G. Bellemare.
DeepMDP: learning continuous latent space models for representation learning.
\emph{ICML}, 2019.

\bibitem{goel1991} V. Goel. Notationality and the information processing mind.
\emph{Minds and Machines} 1(2):129--166, 1991.

\bibitem{goodman1968} N. Goodman. \emph{Languages of Art: An Approach to a Theory of
Symbols}. Bobbs-Merrill, 1968; second edition Hackett, 1976.

\bibitem{harnad1990} S. Harnad. The symbol grounding problem. \emph{Physica D}
42:335--346, 1990.

\bibitem{luo2019} J. Luo, Y. Cao and R. Barzilay. Neural decipherment via minimum-cost
flow: from Ugaritic to Linear B. \emph{ACL}, 2019.

\bibitem{luo2021} J. Luo, F. Hartmann, E. Santus, Y. Cao and R. Barzilay. Deciphering
undersegmented ancient scripts using phonetic prior. \emph{TACL}, 2021.

\bibitem{oord2017} A. van den Oord, O. Vinyals and K. Kavukcuoglu. Neural discrete
representation learning. \emph{NeurIPS}, 2017.

\bibitem{pierce1997} D. Pierce and B. J. Kuipers. Map learning with uninterpreted sensors
and effectors. \emph{Artificial Intelligence} 92:169--229, 1997.

\bibitem{shai2024} A. S. Shai, S. E. Marzen, L. Teixeira, A. Gietelink Oldenziel and
P. M. Riechers. Transformers represent belief state geometry in their residual stream.
\emph{NeurIPS}, 2024.

\bibitem{shalizi2001} C. R. Shalizi and J. P. Crutchfield. Computational mechanics:
pattern and prediction, structure and simplicity. \emph{Journal of Statistical Physics}
104:817--879, 2001.

\bibitem{sharma2024} P. Sharma, S. Gero, R. Payne, D. F. Gruber, D. Rus, A. Torralba and
J. Andreas. Contextual and combinatorial structure in sperm whale vocalisations.
\emph{Nature Communications} 15:3617, 2024.

\bibitem{steels2015} L. Steels. \emph{The Talking Heads Experiment: Origins of Words and
Meanings}. Language Science Press, 2015.

\bibitem{taniguchi2024} T. Taniguchi. Collective predictive coding hypothesis: symbol
emergence as decentralized Bayesian inference. \emph{Frontiers in Robotics and AI}
11:1353870, 2024.

\bibitem{zhang2021} A. Zhang, R. McAllister, R. Calandra, Y. Gal and S. Levine. Learning
invariant representations for reinforcement learning without reconstruction. \emph{ICLR},
2021.

\end{thebibliography}
